
\subsubsection{3.1 Overview}

Our geometric hypothesis has two parts: (1) SSL-SGD creates sharpness anisotropy in the loss landscape, privileging spurious feature directions; (2) SAM applied during SSL pretraining flattens this anisotropy and improves worst-group accuracy. The methodology follows these parts: first, a diagnostic protocol to measure anisotropy annotation-free; second, a training protocol where SAM replaces SGD during SSL pretraining.

\textbf{Connecting to the key insight:} If the loss landscape is geometrically biased toward spurious features, then the SAM perturbation --- which finds the most loss-increasing direction from the current weights --- will naturally find spurious-feature directions more often. This directional preference is measurable, and SAM's flat-minima objective directly counteracts it.

\subsubsection{3.2 Problem Setup}

Let $f_\theta: \mathcal{X} \to \mathbb{R}^d$ be an SSL-pretrained encoder (ResNet-50 backbone, $d=2048$). The dataset $\mathcal{D}$ contains images from a spurious correlation benchmark with underlying group structure $(y, c) \in \{0,1\}^2$, where $y$ is the core label and $c$ is the spurious attribute, correlated at strength $p \approx 0.95$ in training data. Group labels are used only for evaluation, never during training or pretraining.

Let $\mathcal{L}_{SSL}(\theta)$ be the InfoNCE contrastive loss (NT-Xent for SimCLR, InfoNCE for MoCo-v2, self-distillation for DINO). Worst-group accuracy is:
$$\text{WGA} = \min_{g \in \mathcal{G}} \text{Acc}_g(\theta)$$
where $\mathcal{G} = \{(y,c) : y,c \in \{0,1\}\}$ defines four demographic groups.

\subsubsection{3.3 Sharpness Anisotropy Measurement (Annotation-Free)}

We measure directional sharpness anisotropy using three components:

\textbf{Step 1: Linear probe and spurious direction proxy.}
Given a pretrained encoder $f_\theta$, we train a linear probe $h: \mathbb{R}^d \to \mathbb{R}^K$ on the downstream classification task. The per-sample probe loss $\ell_i = \ell(h(f_\theta(x_i)), y_i)$ serves as an annotation-free proxy for spurious correlation reliance: samples in the top-25\% by loss are likely minority group members whose spurious feature conflicts with their label \citep{ghaznavi2023lfr}. Let $S = \{i : \ell_i > Q_{0.75}(\ell)\}$ denote this spurious proxy set.

\textbf{Rationale:} LFR \citep{ghaznavi2023lfr} validates this proxy on Waterbirds and CelebA for supervised ERM representations.

\textbf{Assumption (unvalidated for SSL):} Transfer of this proxy to SSL representations, where the linear probe loss distribution may differ from supervised ERM, requires explicit validation. Section 4.5 includes proxy precision/recall against held-out group labels as a required validation metric. If the proxy fails in SSL representations, the anisotropy ratio measures nothing meaningful.

\textbf{Step 2: SAM perturbation as directional sharpness probe.}
For a given set of samples $T$, we define the directional SAM loss increase as:
$$\Delta\mathcal{L}(T, \theta) = \mathcal{L}(T, \theta + e_T(\theta)) - \mathcal{L}(T, \theta)$$
where $e_T(\theta) = \rho \cdot \nabla_\theta \mathcal{L}(T, \theta) / \|\nabla_\theta \mathcal{L}(T, \theta)\|$ is the SAM perturbation direction (radius $\rho = 0.05$) computed from the gradient on $T$.

\textbf{Step 3: Anisotropy ratio.}
Let $\Delta\mathcal{L}_{\text{spur}} = \Delta\mathcal{L}(S, \theta)$ and $\Delta\mathcal{L}_{\text{rand}} = \frac{1}{N}\sum_{j=1}^{N}\Delta\mathcal{L}(R_j, \theta)$ where $R_j$ are $N=100$ random subsets of the same size as $S$. The \textbf{sharpness anisotropy ratio} is:
$$\text{AR}(\theta) = \frac{\Delta\mathcal{L}_{\text{spur}}}{\Delta\mathcal{L}_{\text{rand}}}$$

$\text{AR} > 1$ indicates the loss landscape is more sensitive to perturbations along spurious directions than random directions. The gate threshold is $\text{AR} > 1.2$ in $\geq 7/9$ SSL method $\times$ dataset combinations.

\textbf{Step 4: Checkpoint correlation.}
We measure $\text{AR}(\theta_t)$ and $\text{WGA}(\theta_t)$ at epochs $\{50, 100, 150, 200\}$, then compute Pearson $r$ between the two sequences. The hypothesis predicts $r < -0.5$.

\textbf{Statistical limitation:} With $n=4$ checkpoints, two-tailed $p<0.05$ requires $|r|>0.950$ --- an extremely tight bar equivalent to a monotonicity test. We report $r$ values with this constraint acknowledged. Future work should increase checkpoint density (e.g., every 25 epochs) for reliable significance testing.

\subsubsection{3.4 SAM-SSL Training Protocol}

We apply SAM during SSL pretraining as a drop-in optimizer replacement: each batch performs a two-step update --- (1) compute gradient and take a perturbation step to the ''sharpest'' neighboring point, then (2) compute gradient at the perturbed point and apply the base SGD update. This is the standard SAM procedure implemented via the davda54/sam wrapper [davda54, 2021] within the izmailovpavel/spurious\_feature\_learning framework.

\textbf{SAM configuration:} rho=0.05 using davda54/sam wrapper; ASAM variant (rho=0.5, adaptive) included as secondary variant.

\textbf{Design rationale --- why SAM during SSL?} We hypothesize that when the loss landscape has higher curvature along spurious directions, SAM's perturbation is more likely to align with those directions, since the gradient direction (which defines the SAM step) is influenced by the curvature structure. This preferential alignment would cause the base optimizer step to flatten those high-curvature spurious directions --- without any explicit group awareness. This mechanistic hypothesis is tested by RQ3. Note that Gatmiry et al. [2024] establish this simplicity-bias mechanism for supervised cross-entropy; whether InfoNCE loss landscapes exhibit analogous behavior is theoretically open and empirically unresolved.

The InfoNCE landscape differs from supervised cross-entropy: it has uniform distribution pressure (von Mises-Fisher concentration) with multiple attractors rather than a single rank-1 attractor. We hypothesize this prevents SAM from simply promoting rank-1 simplicity bias, instead distributing its flattening effect more broadly.

\subsubsection{3.5 Connection to Existing Methods}

Our method is composable with post-hoc debiasing:
\begin{itemize}
\item \textbf{SAM-SSL + LFR/EVaLS}: After SAM pretraining, apply LFR resampling to further improve WGA.
\item \textbf{SAM-SSL + DFR}: Use last-layer retraining on a balanced subset of SAM-SSL features.
\end{itemize}

The primary contribution is the pretraining-level intervention; composability with post-hoc methods provides an upper-bound experiment for future work.

